\documentclass[11pt,a4paper]{article}
\usepackage[utf8]{inputenc}
\usepackage[vietnamese]{babel}
\usepackage{amsmath,amssymb}
\usepackage{geometry}
\geometry{a4paper, total={180mm,260mm}, left=15mm, right=15mm, top=18mm, bottom=18mm}

\usepackage{tikz}
\usetikzlibrary{shapes.geometric, arrows.meta, positioning, calc}
\tikzset{
  arr/.style={-{Stealth[scale=0.8]}},
  every node/.append style={align=center}
}
\usepackage{tcolorbox}
\tcbuselibrary{skins,breakable}
\usepackage{fontawesome}
\usepackage{fancyhdr}
\usepackage{enumitem}

\pagestyle{fancy}
\fancyhf{}
\renewcommand{\headrulewidth}{0.5pt}
\renewcommand{\footrulewidth}{0.5pt}
\fancyhead[L]{\small\bfseries\color{violet!85!black}\faBook\ SỔ TAY KIẾN THỨC TOÁN 2}
\fancyhead[R]{\small\bfseries\color{violet!85!black}Lớp Toán Thầy Lĩnh -- loptoanthaylinh.vn}
\fancyfoot[C]{\small\color{gray}-- Sổ tay kiến thức: Bài 34. Ôn tập hình phẳng --}

\setlength{\parindent}{0pt}
\setlength{\parskip}{0pt}

\begin{document}

% --- BANNER BÀI HỌC ---
\begin{tcolorbox}[
  enhanced,
  colback=violet!8!white,
  colframe=violet!70!black,
  arc=3mm, boxrule=1.2pt,
  left=3mm, right=3mm, top=2mm, bottom=2mm,
  halign=center
]
  {\LARGE\bfseries\color{violet!90!black}\faBookmark\ BÀI 34: ÔN TẬP HÌNH PHẲNG}\\[1mm]
  {\itshape\footnotesize\color{black!75} Ôn thi HK1: Điểm, đoạn thẳng, đường thẳng, đường cong, đường gấp khúc và đếm hình trên 1 trang A4}
\end{tcolorbox}

\vspace{-2mm}

\noindent
\begin{minipage}[t]{0.49\textwidth}\vspace{0pt}
\begin{tcolorbox}[
  enhanced,
  colback=violet!4!white,
  colframe=violet!75!black,
  arc=2mm, boxrule=0.8pt,
  left=2mm, right=2mm, top=1.8mm, bottom=1.8mm,
  title={\footnotesize\bfseries\color{white}\faKey\ 1. Hệ thống các hình phẳng}
]
{\footnotesize
\begin{itemize}\setlength{\itemsep}{1pt}\setlength{\leftskip}{-2.5mm}
  \item \textbf{Điểm:} Kí hiệu bằng chữ cái in hoa ($A, B, C$).
  \item \textbf{Đoạn thẳng:} Có 2 điểm mút, đo được độ dài chính xác.
  \item \textbf{Đường thẳng:} Kéo dài mãi mãi về hai phía.
  \item \textbf{Ba điểm thẳng hàng:} Cùng nằm trên 1 đường thẳng.
  \item \textbf{Đường gấp khúc:} Tổng các đoạn thẳng nối tiếp nhau.
  \item \textbf{Hình tam giác:} Có 3 cạnh, 3 đỉnh.
  \item \textbf{Hình tứ giác:} Có 4 cạnh, 4 đỉnh.
\end{itemize}
}
\end{tcolorbox}
\end{minipage}
\hfill
\begin{minipage}[t]{0.49\textwidth}\vspace{0pt}
\begin{tcolorbox}[
  enhanced,
  colback=teal!4!white,
  colframe=teal!75!black,
  arc=2mm, boxrule=0.8pt,
  left=2mm, right=2mm, top=1.8mm, bottom=1.8mm,
  title={\footnotesize\bfseries\color{white}\faSliders\ 2. Kĩ năng trọng tâm thi HK1}
]
{\footnotesize
\begin{itemize}\setlength{\itemsep}{1pt}\setlength{\leftskip}{-2.5mm}
  \item \textbf{Tính độ dài đường gấp khúc:}\\
  Cộng độ dài tất cả các đoạn thẳng thành phần (nhớ đổi về cùng đơn vị $cm$ hoặc $dm$).
  \item \textbf{Vẽ đoạn thẳng:}\\
  Đặt vạch 0 của thước tại điểm mút đầu tiên.
  \item \textbf{Đếm số hình:}\\
  - Đếm hình tam giác, tứ giác.\\
  - Kĩ thuật đánh số $1, 2, 3$ rồi đếm hình đơn và hình ghép đôi.
\end{itemize}
}
\end{tcolorbox}
\end{minipage}

\vspace{-1.5mm}
\begin{tcolorbox}[
  enhanced,
  colback=orange!4!white,
  colframe=orange!80!black,
  arc=2mm, boxrule=0.8pt,
  left=2mm, right=2mm, top=1.5mm, bottom=1.5mm,
  title={\footnotesize\bfseries\color{white}\faSitemap\ 3. Bảng đối chiếu nhận dạng các yếu tố hình học phẳng HK1}
]
\centering
\begin{tikzpicture}[x=1.5cm, y=0.8cm]
  \node[rectangle, rounded corners=2pt, draw=violet!80!black, fill=violet!15, font=\scriptsize\bfseries] at (0,0) {ĐOẠN THẲNG\\$A \bullet\!\!-\!\!-\!\!-\!\!\bullet B$\\Bị giới hạn 2 đầu};
  \node[rectangle, rounded corners=2pt, draw=teal!80!black, fill=teal!15, font=\scriptsize\bfseries] at (2.5,0) {ĐƯỜNG THẲNG\\$\longleftrightarrow$\\Không giới hạn};
  \node[rectangle, rounded corners=2pt, draw=orange!80!black, fill=orange!15, font=\scriptsize\bfseries] at (5,0) {ĐƯỜNG GẤP KHÚC\\$N$ đoạn thẳng\\Độ dài $=$ Tổng các đoạn};
\end{tikzpicture}
\end{tcolorbox}

\vspace{-1.5mm}
\noindent
\begin{minipage}[t]{0.49\textwidth}\vspace{0pt}
\begin{tcolorbox}[
  enhanced,
  colback=green!4!white,
  colframe=green!60!black,
  arc=2mm, boxrule=0.8pt,
  left=2mm, right=2mm, top=1.8mm, bottom=1.8mm,
  title={\footnotesize\bfseries\color{white}\faCheckCircle\ 4. Ví dụ thực hành chuẩn mực}
]
{\footnotesize
\textbf{1. Tính độ dài đường gấp khúc:}\\
Đường gấp khúc gồm $4$ đoạn thẳng bằng nhau, mỗi đoạn dài $5\text{ cm}$. Tính độ dài đường gấp khúc đó?\\
\textit{Giải:} Độ dài là: $5 + 5 + 5 + 5 = \mathbf{20\text{ cm}}$.
\textbf{2. Đếm hình tứ giác:}\\
Hình vuông có 1 đường thẳng chia đôi:\\
- Gồm $2$ hình chữ nhật nhỏ.\\
- Và $1$ hình vuông lớn ban đầu.\\
$\implies$ Có tất cả $\mathbf{3}$ hình tứ giác!
}
\end{tcolorbox}
\end{minipage}
\hfill
\begin{minipage}[t]{0.49\textwidth}\vspace{0pt}
\begin{tcolorbox}[
  enhanced,
  colback=red!4!white,
  colframe=red!75!black,
  arc=2mm, boxrule=0.8pt,
  left=2mm, right=2mm, top=1.8mm, bottom=1.8mm,
  title={\footnotesize\bfseries\color{white}\faExclamationTriangle\ 5. Mẹo tránh bẫy trắc nghiệm}
]
{\footnotesize
\begin{itemize}\setlength{\itemsep}{1pt}\setlength{\leftskip}{-2.5mm}
  \item \textbf{Hình chữ nhật và hình vuông đều là hình tứ giác!}\\
  Khi đề bài hỏi "có bao nhiêu hình tứ giác", phải đếm cả hình chữ nhật và hình vuông.
  \item Kiểm tra cẩn thận đơn vị đo trước khi cộng tổng độ dài đường gấp khúc.
\end{itemize}
}
\end{tcolorbox}
\end{minipage}


\end{document}
